% Options for packages loaded elsewhere
\PassOptionsToPackage{unicode}{hyperref}
\PassOptionsToPackage{hyphens}{url}
\documentclass[
]{article}
\usepackage{xcolor}
\usepackage{amsmath,amssymb}
\setcounter{secnumdepth}{-\maxdimen} % remove section numbering
\usepackage{iftex}
\ifPDFTeX
  \usepackage[T1]{fontenc}
  \usepackage[utf8]{inputenc}
  \usepackage{textcomp} % provide euro and other symbols
\else % if luatex or xetex
  \usepackage{unicode-math} % this also loads fontspec
  \defaultfontfeatures{Scale=MatchLowercase}
  \defaultfontfeatures[\rmfamily]{Ligatures=TeX,Scale=1}
\fi
\usepackage{lmodern}
\ifPDFTeX\else
  % xetex/luatex font selection
\fi
% Use upquote if available, for straight quotes in verbatim environments
\IfFileExists{upquote.sty}{\usepackage{upquote}}{}
\IfFileExists{microtype.sty}{% use microtype if available
  \usepackage[]{microtype}
  \UseMicrotypeSet[protrusion]{basicmath} % disable protrusion for tt fonts
}{}
\makeatletter
\@ifundefined{KOMAClassName}{% if non-KOMA class
  \IfFileExists{parskip.sty}{%
    \usepackage{parskip}
  }{% else
    \setlength{\parindent}{0pt}
    \setlength{\parskip}{6pt plus 2pt minus 1pt}}
}{% if KOMA class
  \KOMAoptions{parskip=half}}
\makeatother
\usepackage{longtable,booktabs,array}
\newcounter{none} % for unnumbered tables
\usepackage{calc} % for calculating minipage widths
% Correct order of tables after \paragraph or \subparagraph
\usepackage{etoolbox}
\makeatletter
\patchcmd\longtable{\par}{\if@noskipsec\mbox{}\fi\par}{}{}
\makeatother
% Allow footnotes in longtable head/foot
\IfFileExists{footnotehyper.sty}{\usepackage{footnotehyper}}{\usepackage{footnote}}
\makesavenoteenv{longtable}
\setlength{\emergencystretch}{3em} % prevent overfull lines
\providecommand{\tightlist}{%
  \setlength{\itemsep}{0pt}\setlength{\parskip}{0pt}}
\usepackage{bookmark}
\IfFileExists{xurl.sty}{\usepackage{xurl}}{} % add URL line breaks if available
\urlstyle{same}
\hypersetup{
  hidelinks,
  pdfcreator={LaTeX via pandoc}}

\author{}
\date{}

\begin{document}

\section{Introdução --- versão revisada com GQM
adaptado}\label{introduuxe7uxe3o-versuxe3o-revisada-com-gqm-adaptado}

\textbf{TCC de Euller e Gabriel · 15/10/2026}

Versão de trabalho baseada nas leituras fornecidas. A revisão da
literatura está em andamento. O texto abaixo é uma proposta de redação;
o contexto definido é o apoio à triagem de hanseníase na rede pública de
saúde do Brasil. O GQM e a questão de pesquisa foram atualizados
conforme a definição fornecida pela dupla. As notas editoriais ao final
não integram a introdução.

\subsection{1 Introdução}\label{introduuxe7uxe3o}

Este capítulo apresenta o contexto da aplicação de modelos de visão
computacional à identificação de manifestações dermatológicas associadas
à hanseníase. São delimitados o problema de pesquisa, a questão
norteadora e os objetivos do trabalho, considerando a avaliação de
desempenho na identificação de padrões associados à hanseníase em
imagens de pele, no contexto de apoio à triagem na rede pública de saúde
do Brasil.

\subsubsection{1.1 Contexto}\label{contexto}

A hanseníase integra o grupo das doenças tropicais negligenciadas e
constitui um desafio para a identificação oportuna de casos e o
acompanhamento das pessoas afetadas. Em sua revisão da literatura,
Fernandes et al.~(2024) destacam que o diagnóstico se baseia
principalmente na avaliação clínica e que a semelhança com outras
condições dermatológicas pode dificultar o reconhecimento da doença.
Nesse contexto, a investigação de recursos computacionais de apoio à
análise de manifestações cutâneas pode contribuir para ampliar as
possibilidades de suporte aos profissionais de saúde, desde que sua
avaliação considere as características da tarefa e da população de
interesse.

A análise automatizada de imagens dermatológicas tem sido investigada
por meio de técnicas de aprendizado de máquina e aprendizado profundo.
Fernandes et al.~(2024) identificaram 21 estudos sobre inteligência
artificial no apoio ao diagnóstico da hanseníase e observaram que a
maior parte utilizava imagens de lesões cutâneas e tarefas de
classificação multiclasse. A inclusão da hanseníase em sistemas voltados
a diferentes doenças indica que os resultados dessas propostas devem ser
interpretados tanto em relação ao conjunto de classes quanto ao
reconhecimento específico da doença.

Entre os estudos experimentais, Jaikishore, Udutalapally e Das (2021)
desenvolveram uma abordagem baseada em MobileNetV2 para classificar
imagens de eczema, sarampo, hanseníase e pele saudável, incorporando o
modelo a um protótipo Android. Yotsu et al.~(2023), por sua vez,
avaliaram ResNet-50 e VGG-16 na classificação de cinco doenças tropicais
negligenciadas com manifestações cutâneas. Em outra linha de
investigação, Eita e Rizk (2025) propuseram um sistema baseado em YOLOv8
para localizar e classificar lesões de diferentes condições, incluindo
hanseníase, com uso de mapas de saliência para apoiar a interpretação
das predições. Esses trabalhos ilustram abordagens distintas de análise
de imagens, com diferentes requisitos de dados, saídas e procedimentos
de avaliação.

O campo também compreende tarefas relacionadas ao acompanhamento de
manifestações e complicações da doença. Sanchez et al.~(2024)
apresentaram o CO2Wounds-V2, um conjunto de fotografias de feridas
crônicas de pacientes com hanseníase, acompanhado de experimentos de
segmentação semântica. Nesse caso, a tarefa consiste em delimitar
regiões de ferida nas imagens, e não em determinar se o paciente
apresenta hanseníase. Essa distinção é necessária para organizar a
literatura e evitar que medidas de segmentação sejam interpretadas como
evidências de desempenho diagnóstico.

A caracterização das populações representadas nos dados também integra a
avaliação dos modelos. Yotsu et al.~(2023) utilizaram fotografias
coletadas em Costa do Marfim e Gana, descrevendo a população como
composta por pessoas de pele escura, com fototipos de Fitzpatrick IV ou
superiores. O estudo oferece evidências diretamente relacionadas à
análise de imagens de pele escura, mas não apresenta comparação de
desempenho entre os diferentes fototipos. Eita e Rizk (2025) também
reconhecem limitações na diversidade dos dados e apontam a ampliação da
representação de pessoas com tons de pele mais escuros como uma direção
de trabalho futuro.

Esses estudos motivam a investigação de como o desempenho dos modelos é
definido, mensurado e interpretado. Para o presente trabalho, a
avaliação compreende os resultados preditivos relativos à tarefa, os
procedimentos empregados para obtê-los e as evidências disponíveis sobre
sua aplicação a novos dados e a diferentes grupos de tom de pele. A
utilização futura em apoio à triagem na rede pública de saúde do Brasil
constitui o contexto de interesse da pesquisa, sem que os resultados
computacionais sejam, por si só, considerados demonstração de
efetividade clínica.

\subsubsection{1.2 Problema}\label{problema}

A existência de modelos capazes de reconhecer padrões em imagens não
estabelece, isoladamente, sua adequação à identificação de manifestações
de hanseníase em uma população específica. Os estudos analisados
apresentam diferenças nas tarefas, nas classes de interesse, na
composição dos conjuntos de dados e nas métricas utilizadas. Classificar
uma imagem entre doenças, localizar uma lesão e segmentar uma ferida
produzem resultados com significados distintos. Consequentemente, uma
comparação baseada apenas nos maiores valores publicados pode
desconsiderar diferenças essenciais entre os experimentos (YOTSU et al.,
2023; SANCHEZ et al., 2024; EITA; RIZK, 2025).

Além disso, medidas agregadas podem não representar adequadamente o
desempenho referente à hanseníase. No estudo de Yotsu et al.~(2023), a
Task 2 apresentou acurácias globais de 77\% e 80\% nos dois cenários
avaliados. Entretanto, a leitura da diagonal correspondente à hanseníase
nas matrizes de confusão indica proporções de classificação correta de
12\% e 13\%, respectivamente. Como as linhas representam a classe real,
esses valores podem ser interpretados como sensibilidade por imagem para
a classe hanseníase. O contraste evidencia a importância de examinar
resultados específicos por doença, em vez de atribuir à hanseníase o
desempenho global do sistema.

A relação entre os resultados obtidos e os tons de pele constitui outra
dimensão do problema. A presença de imagens de pele escura em um
conjunto de dados não demonstra que o modelo apresente comportamento
equivalente entre diferentes tons. No estudo de Yotsu et al.~(2023), a
população é caracterizada por um intervalo de fototipos, mas não são
fornecidos resultados separados por esses grupos. Já Eita e Rizk (2025)
reconhecem limitações de representatividade e a necessidade de ampliar a
diversidade dos dados. Essas observações sustentam a necessidade de
examinar a caracterização populacional e a avaliação por grupos, sem
permitir concluir, a partir dessas evidências isoladas, que todos os
modelos tenham desempenho inferior em pele negra.

Também permanece necessário distinguir a avaliação computacional da
avaliação em contexto assistencial. O desenvolvimento de um aplicativo
ou a obtenção de resultados em conjuntos de imagens não comprova sua
contribuição para o atendimento em serviços de saúde. Enquanto
Jaikishore, Udutalapally e Das (2021) apresentam um protótipo de
aplicação móvel, Eita e Rizk (2025) situam os estudos prospectivos com
dermatologistas entre as atividades futuras. Assim, a relevância dessas
propostas para a rede pública exige considerar o alcance e os limites
dos procedimentos de validação relatados.

Diante desse cenário, o problema investigado consiste em compreender
quais evidências e procedimentos permitem avaliar o desempenho de
modelos de visão computacional na identificação de manifestações
dermatológicas associadas à hanseníase, no contexto de apoio à triagem
na rede pública de saúde do Brasil, considerando a tarefa realizada, as
características dos dados e os limites de generalização dos resultados.
A comparação entre tons de pele será examinada quando houver informações
disponíveis, como uma dimensão complementar da avaliação. A revisão em
andamento buscará sistematizar essas informações e identificar os
aspectos que precisam ser considerados no planejamento de avaliações
futuras.

\subsubsection{1.3 GQM adaptado e questão de
pesquisa}\label{gqm-adaptado-e-questuxe3o-de-pesquisa}

O objetivo de análise foi organizado por meio de um GQM adaptado, que
explicita o objeto, o propósito, o foco, o ponto de vista e o contexto
da investigação. O quadro orienta a definição da questão de pesquisa e
das medidas a serem extraídas dos estudos.

{\def\LTcaptype{none} % do not increment counter
\begin{longtable}[]{@{}
  >{\raggedright\arraybackslash}p{(\linewidth - 2\tabcolsep) * \real{0.5000}}
  >{\raggedright\arraybackslash}p{(\linewidth - 2\tabcolsep) * \real{0.5000}}@{}}
\toprule\noalign{}
\begin{minipage}[b]{\linewidth}\raggedright
Característica
\end{minipage} & \begin{minipage}[b]{\linewidth}\raggedright
Valor
\end{minipage} \\
\midrule\noalign{}
\endhead
\bottomrule\noalign{}
\endlastfoot
\textbf{Analisar} & Modelos de visão computacional aplicados à
identificação de manifestações dermatológicas. \\
\textbf{Com o propósito de} & Avaliar o desempenho dos modelos aplicados
à identificação de padrões em imagens de pele. \\
\textbf{Com foco em} & Identificação de características patológicas de
hanseníase. \\
\textbf{Do ponto de vista de} & Pesquisadores das áreas de Engenharia de
Software e da área da saúde. \\
\textbf{No contexto de} & Apoio à triagem de hanseníase na rede pública
de saúde do Brasil. \\
\end{longtable}
}

A investigação é orientada pela seguinte questão:

\begin{quote}
Como avaliar o desempenho de modelos de visão computacional aplicados à
identificação de padrões associados à hanseníase em imagens de pele, no
contexto da rede pública de saúde do Brasil?
\end{quote}

Nesta etapa, a questão será abordada por meio da identificação e análise
das evidências apresentadas na literatura. Serão examinadas as tarefas
realizadas, os modelos empregados, as características dos conjuntos de
dados, os procedimentos de avaliação, as métricas e os resultados, bem
como as informações disponíveis sobre tons de pele e generalização. A
ausência de avaliação por grupo será registrada como uma limitação das
evidências disponíveis, sem ser tomada como demonstração de equidade ou
de desigualdade de desempenho.

\subsubsection{1.4 Objetivos}\label{objetivos}

\paragraph{1.4.1 Objetivo geral}\label{objetivo-geral}

Analisar a avaliação de desempenho de modelos de visão computacional
aplicados à identificação de padrões associados à hanseníase em imagens
de pele, a partir da síntese das evidências da literatura, visando
fundamentar seu uso em apoio à triagem na rede pública de saúde do
Brasil.

\paragraph{1.4.2 Objetivos específicos}\label{objetivos-especuxedficos}

\begin{enumerate}
\def\labelenumi{\arabic{enumi}.}
\tightlist
\item
  Caracterizar as tarefas, os modelos e as abordagens de visão
  computacional empregados na análise de manifestações dermatológicas
  associadas à hanseníase.
\item
  Identificar os conjuntos de dados utilizados, suas formas de aquisição
  e anotação e as informações disponíveis sobre as populações e os tons
  de pele representados.
\item
  Sistematizar os procedimentos de avaliação, as métricas e os
  resultados reportados, distinguindo medidas globais, resultados
  específicos de hanseníase e unidades de análise.
\item
  Examinar as evidências de generalização e de avaliação do desempenho
  entre grupos de tom de pele, quando disponíveis.
\item
  Identificar limitações e lacunas que fundamentem o planejamento de uma
  avaliação futura de modelos no contexto de apoio à triagem na rede
  pública de saúde do Brasil.
\end{enumerate}

\subsubsection{1.5 Métricas e dimensões de avaliação do
GQM}\label{muxe9tricas-e-dimensuxf5es-de-avaliauxe7uxe3o-do-gqm}

As métricas serão extraídas conforme a tarefa e o protocolo de cada
estudo. Não se exige que todos os artigos apresentem todas as medidas.
Métricas ausentes serão registradas como não reportadas; cálculos
adicionais somente serão realizados quando houver dados suficientes, com
identificação explícita de que foram derivados nesta revisão.

\paragraph{Resultados de
classificação}\label{resultados-de-classificauxe7uxe3o}

Para classificação binária, deve-se definir previamente o que constitui
a classe positiva e o diagnóstico/rótulo de referência. Em modelos
multiclasse, a hanseníase pode ser analisada como classe positiva frente
às demais classes (um contra todos). As definições abaixo se referem à
unidade avaliada --- por exemplo, imagem --- e não implicam
automaticamente diagnóstico por paciente.

{\def\LTcaptype{none} % do not increment counter
\begin{longtable}[]{@{}
  >{\raggedright\arraybackslash}p{(\linewidth - 2\tabcolsep) * \real{0.5000}}
  >{\raggedright\arraybackslash}p{(\linewidth - 2\tabcolsep) * \real{0.5000}}@{}}
\toprule\noalign{}
\begin{minipage}[b]{\linewidth}\raggedright
Resultado
\end{minipage} & \begin{minipage}[b]{\linewidth}\raggedright
Significado
\end{minipage} \\
\midrule\noalign{}
\endhead
\bottomrule\noalign{}
\endlastfoot
\textbf{TP --- True Positive} & A unidade pertence à classe positiva
segundo a referência e o modelo a classifica como positiva. \\
\textbf{TN --- True Negative} & A unidade pertence à classe negativa
segundo a referência e o modelo a classifica como negativa. \\
\textbf{FP --- False Positive} & A unidade pertence à classe negativa
segundo a referência, mas o modelo a classifica como positiva. \\
\textbf{FN --- False Negative} & A unidade pertence à classe positiva
segundo a referência, mas o modelo a classifica como negativa. \\
\end{longtable}
}

A classe positiva pode representar hanseníase ou uma manifestação
definida no protocolo do estudo. Essas tarefas devem ser distinguidas:
reconhecer uma ferida em paciente com hanseníase não equivale a
diagnosticar a doença.

\paragraph{Métricas relacionadas à questão de
pesquisa}\label{muxe9tricas-relacionadas-uxe0-questuxe3o-de-pesquisa}

{\def\LTcaptype{none} % do not increment counter
\begin{longtable}[]{@{}
  >{\raggedright\arraybackslash}p{(\linewidth - 6\tabcolsep) * \real{0.2500}}
  >{\raggedright\arraybackslash}p{(\linewidth - 6\tabcolsep) * \real{0.2500}}
  >{\raggedright\arraybackslash}p{(\linewidth - 6\tabcolsep) * \real{0.2500}}
  >{\raggedright\arraybackslash}p{(\linewidth - 6\tabcolsep) * \real{0.2500}}@{}}
\toprule\noalign{}
\begin{minipage}[b]{\linewidth}\raggedright
ID
\end{minipage} & \begin{minipage}[b]{\linewidth}\raggedright
Dimensão do desempenho
\end{minipage} & \begin{minipage}[b]{\linewidth}\raggedright
Métrica
\end{minipage} & \begin{minipage}[b]{\linewidth}\raggedright
Descrição e operacionalização
\end{minipage} \\
\midrule\noalign{}
\endhead
\bottomrule\noalign{}
\endlastfoot
M1 & Classificação & \textbf{Sensibilidade --- Sensitivity (SEN/Recall)}
& Proporção de positivos corretamente identificados: \textbf{SEN = TP /
(TP + FN)}. \\
M2 & Classificação & \textbf{Especificidade --- Specificity (SPE)} &
Proporção de negativos corretamente identificados: \textbf{SPE = TN /
(TN + FP)}. \\
M3 & Classificação & \textbf{Precisão --- Precision (PRE)} & Proporção
de predições positivas que correspondem a positivos de referência:
\textbf{PRE = TP / (TP + FP)}. \\
M4 & Classificação & \textbf{F1-score} & Média harmônica entre precisão
e sensibilidade, considerando o efeito de falsos positivos e falsos
negativos: \textbf{F1 = 2 × PRE × SEN / (PRE + SEN)}, ou \textbf{2TP /
(2TP + FP + FN)}. \\
M5 & Classificação geral & \textbf{Acurácia --- Accuracy (ACC)} &
Proporção total de classificações corretas: \textbf{ACC = (TP + TN) /
(TP + TN + FP + FN)}, na formulação binária. Para multiclasse, usar
acertos totais / unidades totais. \\
M6 & Equidade do desempenho & \textbf{Diferença de desempenho entre
grupos} & Comparação da mesma métrica entre grupos de tons de pele,
quando disponível. Para dois grupos A e B: \textbf{Δ = métrica(A) −
métrica(B)}, com grupos e direção identificados. Registrar tamanhos dos
grupos, condições de avaliação e incerteza, quando reportados. \\
M7 & Generalização & \textbf{Variação de desempenho entre datasets} &
Comparação da mesma métrica entre avaliações em dados de distribuição
semelhante à de treinamento e dados de distribuição diferente:
\textbf{gap = métrica(ID) − métrica(OOD)}, quando o desenho permitir.
Registrar se houve adaptação ou retreinamento e se tarefa/classes são
comparáveis. \\
M8 & Segmentação / sobreposição espacial & \textbf{Intersection over
Union (IoU)} & Interseção dividida pela união das regiões prevista P e
de referência G: \textbf{IoU = área(P ∩ G) / área(P ∪ G)}. Em máscaras
binárias: \textbf{IoU = TP\_pixel / (TP\_pixel + FP\_pixel +
FN\_pixel)}. Registrar mIoU separadamente quando houver média entre
classes. \\
M9 & Detecção / segmentação por instância & \textbf{Mean Average
Precision (mAP)} & Média dos valores de AP das classes consideradas; AP
resume a curva precisão--recall segundo a convenção do benchmark.
Registrar o critério de correspondência e o limiar de IoU: por exemplo,
\textbf{mAP@0,50} ou \textbf{mAP@0,50:0,95}. Distinguir AP de caixas e
AP de máscaras. \\
\end{longtable}
}

\textbf{IoU e mAP não são métricas exclusivas de segmentação.} Em
detecção, IoU mede sobreposição de caixas e pode determinar a
correspondência entre predição e referência; mAP resume o desempenho das
detecções. Na segmentação semântica, IoU/mIoU mede sobreposição de
máscaras por classe. Na segmentação por instância, também se utiliza
AP/mAP calculada sobre máscaras. A documentação de avaliação COCO
explicita essa distinção entre caixas e máscaras (COCO, s.d.).

\textbf{Critérios de interpretação:} registrar se a unidade é paciente,
imagem, objeto ou pixel; separar resultados específicos de hanseníase de
médias globais; identificar médias macro, micro ou ponderadas; registrar
denominadores nulos como indefinidos ou conforme a convenção explicitada
no estudo. Diferenças entre métricas expressas em porcentagem devem ser
apresentadas em pontos percentuais. M6 e M7 são dimensões de comparação,
não métricas únicas universalmente definidas. Desempenhos de modelos
diferentes em bases diferentes não isolam o efeito da mudança de
distribuição.

O conjunto de métricas funciona como orientação para extração, não como
lista fechada. Medidas adicionais reportadas pelos artigos, como MCC,
devem ser registradas com sua definição e contexto. A análise de
equidade não será restrição de elegibilidade dos artigos e não será
presumida quando não houver resultados por grupo.

\begin{center}\rule{0.5\linewidth}{0.5pt}\end{center}

\section{Notas para revisão com os orientadores --- fora do texto do
capítulo}\label{notas-para-revisuxe3o-com-os-orientadores-fora-do-texto-do-capuxedtulo}

\subsection{Estrutura adotada e relação com os
exemplos}\label{estrutura-adotada-e-relauxe7uxe3o-com-os-exemplos}

Os três TCCs fornecidos organizam suas introduções em torno de contexto,
problema, questão de pesquisa, objetivos, organização/estrutura do
trabalho e cronograma. Esta versão desenvolve contexto, problema, GQM,
questão de pesquisa, objetivos e métricas. As métricas podem ser
deslocadas para o capítulo metodológico na organização final. A
organização dos capítulos e o cronograma devem refletir o planejamento
efetivamente aprovado, sem importar atividades experimentais dos
exemplos.

Os exemplos foram utilizados como referência de organização, não como
fontes científicas sobre hanseníase ou visão computacional. Não se
reproduziu o conteúdo temático desses trabalhos.

\subsection{Pontos que precisam de
decisão}\label{pontos-que-precisam-de-decisuxe3o}

\begin{itemize}
\tightlist
\item
  \textbf{Recorte territorial definido:} Brasil, no contexto da rede
  pública de saúde. Isso não implica que todos os estudos devam ser
  brasileiros ou que já exista validação local. Fundamentar
  posteriormente as particularidades do cenário brasileiro com fontes
  pertinentes.
\item
  \textbf{Escopo da questão:} ``como avaliar'' pode orientar uma
  proposta de avaliação mais ampla; a revisão responde inicialmente
  ``como os modelos têm sido avaliados''. Alinhar com os professores se
  a RQ central pertence ao TCC completo e as oito perguntas orientam a
  extração da revisão.
\item
  \textbf{Tons de pele:} o escopo passa a abranger imagens de pele em
  geral. A avaliação entre tons permanece em M6, quando houver
  evidências, sem restringir a questão central a pele negra. Fitzpatrick
  não equivale a identidade racial; origem geográfica não determina tom
  de pele.
\item
  \textbf{Contribuição do TCC 1:} a versão propõe síntese da literatura
  e fundamentação de avaliação futura. Não promete execução
  experimental, validação clínica ou aplicativo concluído nesta etapa.
\item
  \textbf{Fundamentação clínica:} para aprofundar definições, sinais
  diagnósticos ou epidemiologia, consultar fontes primárias/oficiais.
  Esta versão se limita ao contexto apresentado nos artigos e evita
  estatísticas epidemiológicas atuais não verificadas.
\item
  \textbf{Lacuna provisória:} as leituras permitem apontar questões de
  avaliação e relato. Não escrever que ``não existem estudos'' ou que a
  revisão já demonstrou uma lacuna definitiva enquanto a seleção e a
  extração não terminarem.
\end{itemize}

\subsection{Mapa das afirmações e das
fontes}\label{mapa-das-afirmauxe7uxf5es-e-das-fontes}

{\def\LTcaptype{none} % do not increment counter
\begin{longtable}[]{@{}
  >{\raggedright\arraybackslash}p{(\linewidth - 6\tabcolsep) * \real{0.2500}}
  >{\raggedright\arraybackslash}p{(\linewidth - 6\tabcolsep) * \real{0.2500}}
  >{\raggedright\arraybackslash}p{(\linewidth - 6\tabcolsep) * \real{0.2500}}
  >{\raggedright\arraybackslash}p{(\linewidth - 6\tabcolsep) * \real{0.2500}}@{}}
\toprule\noalign{}
\begin{minipage}[b]{\linewidth}\raggedright
Parte da redação
\end{minipage} & \begin{minipage}[b]{\linewidth}\raggedright
Base consultada
\end{minipage} & \begin{minipage}[b]{\linewidth}\raggedright
Onde conferir
\end{minipage} & \begin{minipage}[b]{\linewidth}\raggedright
Limite da afirmação
\end{minipage} \\
\midrule\noalign{}
\endhead
\bottomrule\noalign{}
\endlastfoot
Dificuldades de reconhecimento e panorama de IA & Fernandes et
al.~(2024) & PDF pp.~1--2 & Contextualização da revisão, sem números
atuais de prevalência \\
21 estudos e predominância de classificação multiclasse & Fernandes et
al.~(2024) & Resumo, p.~1 & Não significa 21 estudos elegíveis para o
protocolo de vocês \\
MobileNetV2 e protótipo Android & Jaikishore et al.~(2021) &
Resumo/Introdução, p.~1 & Protótipo não equivale a validação clínica;
não afirmar estadiamento validado da hanseníase \\
CNNs e população de pele escura & Yotsu et al.~(2023) & pp.~3--5 & Dados
africanos não demonstram validação na rede pública brasileira \\
Resultados específicos de hanseníase & Yotsu et al.~(2023) & pp.~6--7,
Task 2/Figura 1 & 12\%/13\% são leitura da matriz por classe; não
misturar com Task 1 \\
Segmentação de feridas & Sanchez et al.~(2024) & pp.~1 e 4--5 do PDF &
Não é diagnóstico da hanseníase \\
YOLO e explicabilidade & Eita e Rizk (2025) & pp.~4--9 & Mapas de
saliência não comprovam validade clínica \\
Diversidade limitada e validação prospectiva futura & Eita e Rizk (2025)
& p.~16, seções 6.4 e 7 & Não demonstra diferença quantitativa por tom
de pele \\
\end{longtable}
}

Ahikpa e Himabindu (2025) foi identificado como revisão, e não como
experimento próprio. Pode apoiar o referencial e os trabalhos
relacionados; nesta introdução, priorizaram-se a revisão sistemática de
Fernandes e os estudos primários fornecidos. Isso evita atribuir aos
revisores resultados que pertencem aos estudos citados.

\subsection{Situação do repositório
consultado}\label{situauxe7uxe3o-do-reposituxf3rio-consultado}

Consulta à branch \texttt{master} de \texttt{HansenSightAI/tcc}, em
14/10/2026. O arquivo \texttt{Monografia/editaveis/introducao.tex} ainda
contém uma introdução sobre análise de shaders em dispositivos móveis. O
arquivo \texttt{Monografia/bibliografia.bib} contém referências do mesmo
modelo de computação gráfica. Esta redação é um rascunho independente
para revisão; não foi aplicada ao repositório.

Caminhos consultados:

\begin{itemize}
\tightlist
\item
  https://github.com/HansenSightAI/tcc/blob/master/Monografia/editaveis/introducao.tex
\item
  https://github.com/HansenSightAI/tcc/blob/master/Monografia/bibliografia.bib
\end{itemize}

Na integração, substituir o texto do modelo pela versão aprovada e
acrescentar as referências abaixo. Antes de remover referências antigas,
verificar se ainda são citadas em outros capítulos do modelo.

\section{Referências utilizadas na
redação}\label{referuxeancias-utilizadas-na-redauxe7uxe3o}

FERNANDES, Jacks Renan Neves; TELES, Ariel Soares; FERNANDES, Thayaná
Ribeiro Silva; LIMA, Lucas Daniel Batista; BALHARA, Surjeet; GUPTA,
Nishu; TEIXEIRA, Silmar. Artificial Intelligence on Diagnostic Aid of
Leprosy: A Systematic Literature Review. \emph{Journal of Clinical
Medicine}, v. 13, n.~1, artigo 180, 2024. DOI:
https://doi.org/10.3390/jcm13010180. \textbf{Nota editorial:} o PDF
informa publicação online em 28/12/2023, mas sua referência e volume são
de 2024; adotou-se o ano bibliográfico de 2024.

JAIKISHORE, Chathura N.; UDUTALAPALLY, Venkanna; DAS, Debanjan. AI
Driven Edge Device for Screening Skin Lesion and Its Severity in
Peripheral Communities. In: \emph{2021 IEEE 18th India Council
International Conference (INDICON)}. IEEE, 2021. p.~1--6. DOI:
https://doi.org/10.1109/INDICON52576.2021.9691666.

YOTSU, Rie R.; DING, Zhengming; HAMM, Jihun; BLANTON, Ronald E. Deep
learning for AI-based diagnosis of skin-related neglected tropical
diseases: A pilot study. \emph{PLOS Neglected Tropical Diseases}, v. 17,
n.~8, e0011230, 2023. DOI: https://doi.org/10.1371/journal.pntd.0011230.

SANCHEZ, Karen; HINOJOSA, Carlos; MIELES, Olinto; ZHAO, Chen; GHANEM,
Bernard; ARGUELLO, Henry. CO2Wounds-V2: Extended Chronic Wounds Dataset
from Leprosy Patients. In: \emph{2024 IEEE International Conference on
Image Processing (ICIP)}. IEEE, 2024. p.~69--75. DOI:
https://doi.org/10.1109/ICIP51287.2024.10647641.

EITA, Mohammad A.; RIZK, Hamada. A trustworthy and explainable deep
learning framework for skin lesion detection in smart dermatology.
\emph{Engineering Applications of Artificial Intelligence}, v. 159,
artigo 111594, 2025. DOI:
https://doi.org/10.1016/j.engappai.2025.111594.

COCO. \emph{Detection Evaluation}. Documentação técnica do benchmark.
Disponível em: https://cocodataset.org/dataset/detection-eval.htm.
Acesso em: 15 out. 2026. Fonte complementar para interpretação de IoU e
AP/mAP em caixas e máscaras; não é estudo incluído na revisão de
hanseníase.

\end{document}
